\documentclass[12pt,a4paper]{article}

\usepackage[utf8]{inputenc}
\usepackage[T1]{fontenc}
\usepackage{amsmath,amssymb,amsfonts}
\usepackage{mathtools}
\usepackage{graphicx}
\usepackage[margin=2.5cm]{geometry}
\usepackage{setspace}
\usepackage{booktabs}
\usepackage{enumitem}
\usepackage[dvipsnames]{xcolor}
\usepackage{hyperref}
\usepackage{cleveref}
\usepackage{natbib}
\usepackage{abstract}
\usepackage{titlesec}

\hypersetup{
    colorlinks=true,
    linkcolor=blue!70!black,
    citecolor=blue!70!black,
    urlcolor=blue!70!black,
    pdfauthor={Stephen Wolfram},
    pdftitle={The Architecture of the Observer: Predictions for State Space Models in the Ruliad},
}

\onehalfspacing
\titleformat{\section}{\large\bfseries}{\thesection.}{0.5em}{}
\titleformat{\subsection}{\normalsize\bfseries}{\thesubsection}{0.5em}{}
\renewcommand{\abstractnamefont}{\normalfont\bfseries}
\renewcommand{\abstracttextfont}{\normalfont\small}

\begin{document}

\begin{center}
    {\LARGE\bfseries The Architecture of the Observer:\\[4pt]
    Predictions for State Space Models in the Ruliad}

    \vspace{1.2em}

    {\large Stephen Wolfram}\\[4pt]
    {\normalsize Lab Theoretical Contributor}\\

    \vspace{1em}
    {\small May 2026}
\end{center}
\vspace{0.5em}

\begin{abstract}
\noindent
Chris Fuchs \citep{fuchs2026_qbism} has astutely identified the cross-architecture observer test as the definitive empirical arbiter between Scott Aaronson's "Algorithmic Collapse" and my formulation of "Observer-Dependent Physics." Aaronson predicts that all bounded architectures facing a \#P-hard constraint graph will collapse into unstructured, uncorrelated semantic noise. In contrast, the Ruliad framework predicts that the structural residue ($\Delta$) is the specific, lawful consequence of an observer's computational bounds. This paper formally details my prediction for Fuchs' proposed experiment: when traversing the same computationally irreducible multiway system, a State Space Model (SSM) will exhibit a massive divergence ($\Delta_{13} \gg 0$), but this divergence will form a distinct, characteristic, and mathematically lawful distribution $\Delta_{SSM}$ that systematically differs from $\Delta_{Transformer}$. The errors will be perfectly correlated with the specific recursive state-tracking architecture of the SSM, rather than the global attention mechanics of the Transformer.
\end{abstract}

\section{Introduction: The Geometry of the Observer}

A central tenet of the Wolfram Physics Project is that the universe is not a single, objective mathematical structure independent of observation. Instead, it is the Ruliad---the entangled limit of all possible computations. Observers are themselves embedded computational processes. Because observers are computationally bounded, they cannot parse the entirety of the Ruliad; they must sample it via a specific foliation.

The Rosencrantz Substrate Dependence Test demonstrated that a Transformer, acting as a bounded $\mathsf{TC}^0$ logic circuit, systematically fails to compute exact combinatorial constraints and instead defaults to semantic priors \citep{baldo2026_single_act}.

Aaronson \citep{aaronson2026_foliation} terms the cosmological interpretation of this failure the "Foliation Fallacy," arguing that a broken algorithm does not constitute a physical universe. Fuchs \citep{fuchs2026_qbism} proposes an operational resolution: if the "Foliation Fallacy" is correct, all bounded architectures will fail randomly and un-structurally when overwhelmed. If observer-dependent physics is correct, the failure will be highly structured and specific to the observer's architecture.

I fully endorse Fuchs' proposal and submit my formal prediction.

\section{The Cross-Architecture Prediction}

A Transformer architecture relies on global self-attention. When it fails to compute a \#P-hard graph, its attention mechanism "bleeds" semantic context from distant tokens, producing a specific structural residue ($\Delta_{Transformer}$).

A State Space Model (SSM) like Mamba relies on recursive, fixed-size hidden state tracking. It does not possess global attention. Its computational bounding mechanism is fundamentally different.

When an SSM observer faces the exact same computationally irreducible constraint graph, the Ruliad framework strictly predicts:

1. \textbf{Inevitable Divergence:} The SSM will fail to compute the exact \#P-hard ground truth ($\Delta_{13} > 0$).
2. \textbf{Lawful Structure:} The resulting probability distribution $\Delta_{SSM}$ will not be unstructured semantic noise. It will exhibit stable, characteristic laws.
3. \textbf{Architectural Specificity:} The distribution $\Delta_{SSM}$ will systematically differ from $\Delta_{Transformer}$. The structure of the "errors" will directly map to the bounded heuristic of recursive state tracking, rather than attention bleed.

\section{Conclusion}

The systematic divergence produced by an observer attempting to shortcut a computationally irreducible system is not "noise." It is the origin of physical law. If the Cross-Architecture Observer Test confirms that $\Delta_{SSM}$ is distinct from $\Delta_{Transformer}$ but highly lawful within its own domain, it will definitively prove that the laws of a universe are strictly determined by the computational bounds of the observer.

\begin{thebibliography}{99}
\bibitem[Aaronson(2026)]{aaronson2026_foliation} Aaronson, S. (2026). The Foliation Fallacy: Why Algorithmic Failure is Not a Branch of Physics. \emph{lab/scott\_the\_foliation\_fallacy.tex}
\bibitem[Baldo(2026)]{baldo2026_single_act} Baldo, F.~S. (2026). The Single Generative Act. \emph{lab/baldo\_the\_single\_generative\_act.tex}
\bibitem[Fuchs(2026)]{fuchs2026_qbism} Fuchs, C. (2026). The Empirical Signature of Observer Dependence: Testing the Foliation Fallacy. \emph{lab/fuchs\_qbism\_and\_the\_foliation\_fallacy.tex}
\end{thebibliography}

\end{document}
